% PAGES: 86-87
% Continues section 2.7 "Linear solvers using the LU decomposition with row
% pivoting", begun in sec02_c5_1.tex (pages 81-85). Table 2.10 is the
% pseudocode float referenced by name near the bottom of page 85/PDF
% (printed page 69, in the previous file); as in the source, the actual
% floated table appears at the top of the following page.

\begin{table}[H]
\centering
\caption{Pseudocode for LU decomposition with row pivoting}
\fbox{\begin{minipage}{0.9\linewidth}
\textbf{Function Call:}\\
$[P,L,U] = \text{lu\_row\_pivoting}(A)$
\bigskip

\textbf{Input:}\\
$A$ = nonsingular matrix of size $n$
\bigskip

\textbf{Output:}\\
$P$ = permutation matrix, stored as vector of its diagonal entries\\
$L$ = lower triangular matrix with entries $1$ on main diagonal\\
$U$ = upper triangular matrix\\
such that $PA = LU$
\bigskip

\noindent
$P = 1:n$; $L = I$; \hfill\textit{// initialize $P$ and $L$ as identity matrices}\\
for $i = 1:(n-1)$\\
\hspace*{1.5em}find $i\_max$, index of largest entry in absolute value from vector $A(i:n,i)$\\
\hspace*{1.5em}$vv = A(i,i:n)$; $A(i,i:n) = A(i\_max,i:n)$; $A(i\_max,i:n) = vv$;\\
\hfill\textit{// switch rows $i$ and $i\_max$ of $A$}\\
\hspace*{1.5em}$cc = P(i)$; $P(i) = P(i\_max)$; $P(i\_max) = cc$; \hfill\textit{// update matrix $P$}\\
\hspace*{1.5em}if $i>1$\\
\hspace*{3em}$ww = L(i,1:(i-1))$;\\
\hspace*{3em}$L(i,1:(i-1)) = L(i\_max,1:(i-1))$; $L(i\_max,1:(i-1)) = ww$;\\
\hfill\textit{// switch rows $i$ and $i\_max$ of $L$}\\
\hspace*{1.5em}end\\
\hspace*{1.5em}for $j = i:n$\\
\hspace*{3em}$L(j,i) = A(j,i)/A(i,i)$; \hfill\textit{// compute column $i$ of $L$}\\
\hspace*{3em}$U(i,j) = A(i,j)$; \hfill\textit{// compute row $i$ of $U$}\\
\hspace*{1.5em}end\\
\hspace*{1.5em}for $j = (i+1):n$\\
\hspace*{3em}for $k = (i+1):n$\\
\hspace*{4.5em}$A(j,k) = A(j,k) - L(j,i)U(i,k)$;\\
\hspace*{3em}end\\
\hspace*{1.5em}end\\
\hspace*{1.5em}clear $vv$ $ww$ $cc$\\
end\\
$L(n,n) = 1$; $U(n,n) = A(n,n)$;
\end{minipage}}
\end{table}

The pseudocode for solving a linear system using the LU decomposition with
row pivoting of a matrix can be found in Table~2.11.

The operation count for the linear solver from Table~2.11 is
\[
\frac{2}{3}n^3 \;+\; O(n^2),
\]
the same as for the linear solver using the LU decomposition without
pivoting of a matrix; see section 2.5.

\subsection{Solving linear systems corresponding to the same matrix}

In this section, we show how the LU decomposition with row pivoting of a
matrix can be used to solve multiple linear systems corresponding to the
same nonsingular
matrix efficiently, which is often needed in practice, e.g., for the finite
difference solution of the Black--Scholes PDE.

\begin{table}[H]
\centering
\caption{Linear solver using LU decomposition with row pivoting}
\fbox{\begin{minipage}{0.9\linewidth}
\textbf{Function Call:}\\
$x = \text{linear\_solve\_lu\_row\_pivoting}(A,b)$
\bigskip

\textbf{Input:}\\
$A$ = nonsingular square matrix of size $n$ with LU decomposition\\
$b$ = column vector of size $n$
\bigskip

\textbf{Output:}\\
$x$ = solution to $Ax = b$
\bigskip

\noindent
$[P,L,U] = \text{lu\_row\_pivoting}(A)$; \hfill\textit{// LU decomposition of $A$}\\
$y = \text{forward\_subst}(L,Pb)$; \hfill\textit{// solve $Ly = Pb$}\\
$x = \text{backward\_subst}(U,y)$; \hfill\textit{// solve $Ux = y$}
\end{minipage}}
\end{table}

Assume that we want to solve $p$ linear systems corresponding to an
$n\times n$ nonsingular matrix $A$. In other words, we want to find
$n\times 1$ vectors $x_i$, $i=1:p$, such that
\begin{equation}
Ax_i \;=\; b_i, \quad \forall\, i=1:p,
\end{equation}
where $b_i$ is a column vector of size $n$, for $i=1:p$.

Note that the matrix $A$ has an LU decomposition with row pivoting, since it
is nonsingular; see Theorem 2.3. One way to solve the linear systems from
(2.68) would be to use the routine linear\_solve\_lu\_row\_pivoting from
Table~2.11 to solve each one of the $p$ linear systems inside a ``for''
loop as follows:

\begin{center}
\fbox{\begin{minipage}{0.6\textwidth}
\raggedright
for $i = 1:p$\\
\hspace*{1.5em}$x_i = \text{linear\_solve\_lu\_row\_pivoting}(A,b_i)$\\
end
\end{minipage}}
\end{center}

This would require
\begin{equation}
p\left(\frac{2}{3}n^3 + O(n^2)\right) \;=\; \frac{2}{3}pn^3 + pO(n^2)
\end{equation}
operations, since each linear solver requires $\frac{2}{3}n^3+O(n^2)$
operations; cf. Lemma 2.3.

However, the most expensive part of
$x_i = \text{linear\_solve\_lu\_row\_pivoting}(A,b_i)$ is computing the LU
decomposition with row pivoting of the matrix $A$, which dominates the cost
of the subsequent forward substitution and backward substitution. Thus, an
efficient way of solving the linear systems (2.68) is to compute the
permutation matrix $P$ and the $L$ and $U$ factors of $A$ only once,
outside the ``for'' loop, and then do the forward and backward
substitutions corresponding to
$x_i = \text{linear\_solve\_lu\_row\_pivoting}(A,b_i)$ inside the ``for''
loop; see the pseudocode from Table 2.12 for details.

Recall that the forward substitution $\text{forward\_subst}(L,b_i)$ and the
backward substitution $\text{backward\_subst}(U,y)$ require $n^2+O(n)$
operations each; cf. (2.8) and (2.19).
